\section{What a further global conjecture still needs}
\label{sec:global}

\subsection{The cost of missing dimensions}
The source's last comparison uses only $n\ge D/2$ when dividing by the number of chosen vectors. Reading that inequality alone can obscure the stronger hypothesis already used to establish the lower bound. The declaration
\up{grh/src/New/Determinant/Twist/LongForms.lean}{\file{integral_forms_long}}
requires
\begin{equation}
n\ge(1-4/b^2)D.
\label{eq:near-full}
\end{equation}
Each missing basis vector has completion cost
\[
 K_{\mathrm{comp}}=2x_1^2b^2+b.
\]
The source lemma \file{mul_completion_cost_le} combines this expensive cost with the very small rank deficit. It is not enough to transplant a statement saying that a fixed positive fraction of target forms is spanned.

The general determinant principle gives an exact way to track the tradeoff.
\begin{proposition}[Rank deficit versus completion cost]
\label{prop:rank}
Suppose $D=n+r$, $n>0$, $m>0$, and an upper bound and a lower determinant estimate take the forms
\[
 c\le n(A-mR+U),\qquad
 0\le c+2I+rK-DA+V.
\]
Then
\[
 nmR\le nU+2I+V+r(K-A).
\]
If $A\ge0$, this implies
\[
 R\le\frac{U+(2I+V)/n+(r/n)K}{m}.
\]
\end{proposition}
\begin{proof}
Substitute the upper bound for $c$ in the second inequality, expand $D=n+r$, and cancel $nA$. For the second conclusion, discard the nonpositive term $-rA$ and divide by $nm$.
\end{proof}
The proposition is an elementary deduction once the source-compatible estimates are available. Its value is diagnostic: an extension must achieve a small rank deficit, cheaper completion vectors, or a compensating estimate for the negative $rA$ term. Merely improving the torsion selector does not pay a large archimedean completion bill.

\subsection{Only selected twists need injective lifts}
In the generic
\up{grh/src/New/Reduction/Lattices/TwistedLattice.lean}{\file{covolume_add_index_nonneg}},
injectivity of the lift $\Theta_\xi$ is assumed for every twist. The proof uses it only at the trivial twist and at each twist selected for a finite-place embedding. Those are the instances needed to make the reference Gram determinant and selected determinant ratios nonzero.

Thus the proof admits a weaker injectivity hypothesis: it suffices to have injectivity at those selected twists. A family may contain bad twists if the arithmetic selection avoids them. This is a source-derived simplification, not a newly compiled Lean refactor.

The essential compatibility requirements remain:
\begin{itemize}
\item one reference Gram determinant and one basis;
\item index matrices that give integrality for the required local twists;
\item integrality at an embedding and its conjugate;
\item a height bound using a common global multiplier.
\end{itemize}
Independent favorable choices at different primes do not by themselves satisfy this global height statement. The source's \file{HeightLe} explicitly includes the common multiplier.

\subsection{A concrete nonabelian benchmark}
The splitting field of $X^3-X-1$ remains an economical test case. The polynomial has no rational root and has discriminant $-23$, hence its Galois group is $S_3$. Its two-dimensional irreducible representation has character values $2,0,-1$ on the identity, transpositions, and $3$-cycles. The associated degree-two Artin factor is beyond the exported Dirichlet theorem.

The previous reports already reduce the desired zero location to a class-sensitive or Frobenius-sensitive one-sided prime estimate. Repeating that analytic reduction would not prove the estimate. The decisive next object would be a family of integral forms whose local divisibility records this nonabelian information and whose global norms, span, and denominators admit compatible bounds.

For this benchmark the required tasks are now precise:
\begin{enumerate}
\item Define the global lattice, its integral structures, and the statistic recorded by its local divisibility. Congruence information alone cannot be substituted for nonabelian Frobenius data without a proof that it retains the relevant classes.
\item Produce sufficiently many short independent vectors, with an explicit rank-deficit bound and completion cost satisfying Proposition~\ref{prop:rank}.
\item Prove the lower determinant estimate and a uniform total denominator budget. If a fixed coefficient field is available, the present exact specialization theorem can improve this step.
\item Check compatibility of the same twists with divisibility, nonvanishing, paired integrality, and global heights.
\item Apply the already-developed analytic zero-exclusion criterion, retaining its conditions that prevent cancellation of zero contributions.
\end{enumerate}
No such new lattice construction is supplied by this report or by the inspected nonabelian detector reports.

\subsection{Selmer groups and other proposed targets}
The general twisting theorem in
\up{grh/src/New/Congruences/Selmer/GeneralTheorem.lean}{\file{GeneralTheorem.lean}}
has substantial representation-theoretic hypotheses. The actual
\up{grh/src/New/Congruences/Selmer/Line.lean}{\file{TwistedLengthLe}}
bounds a divisible-coefficient Selmer module, not merely one finite torsion level. Finite length therefore rules out a positive divisible corank for the selected object. It does not force its finite torsion subgroup to vanish.

The distinction is concrete: a nonzero residue module over a coefficient ring of ramification index $e'$ has normalized length $1/e'$. Such lengths can tend to zero while the modules remain nonzero. Furthermore, changing the twist or its ramified primes changes the Selmer problem. A comparison to one fixed arithmetic group must be proved before a family of small bounds yields its vanishing.

\begin{center}
\begin{tabular}{@{}>{\raggedright\arraybackslash}p{.25\linewidth}>{\raggedright\arraybackslash}p{.66\linewidth}@{}}
\toprule
Target & Present obstruction\\
\midrule
General Hecke/Artin RH & No compatible global construction giving the new Frobenius-sensitive divisibility and determinant inequality.\\
Modular and symmetric-power RH & The required Hecke-coefficient analogue of the arithmetic comparison is not established.\\
BSD-type finiteness or vanishing & Conventional local conditions must be identified; conjugate self-duality and scalar-inertia hypotheses can exclude the intended representation.\\
Vandiver-type assertions & Parity, exceptional coefficient primes, and changes of field remain; exact local selection alone does not supply the required class-group comparison.\\
Candidate $K$-group vanishing & The strict-Selmer-to-Hecke bridge and a valid finite certificate for the exact arithmetic object remain separate obligations.\\
Twin primes, Goldbach, Elliott--Halberstam & The source's one-variable weighted prime sums do not provide the needed correlations or uniformity.\\
\bottomrule
\end{tabular}
\end{center}

The prior corpus includes further downstream arithmetic classifications and sparse CM results. They are recorded as claims of those reports in Appendix~A. They do not replace the missing global construction, and this report does not recertify them as additional new results.

\subsection{Assessment}
The local question has advanced: the coefficient-field obstruction is exact at zero tolerance, and the multivariable field-extension tower is unnecessary. The major global question is still an arithmetic one. The most useful next step is a complete construction for one small nonabelian target, with all four cost components recorded: divisibility gain, rank deficit, completion cost, and denominator loss. Another analytic detector can be useful only after it changes one of these obligations.

